\begin{lemma}
Under the hypotheses of \(\noderef{HugeFamilySizeBound}\), and assuming also
\(\varepsilon_1\ge0\), suppose the sample parameters are the rounded hierarchy
parameters
\[
m=\noderef{TwoBiteNaturalReducedVertexCount}(n),
\qquad
p=\noderef{TwoBiteEdgeProbability}\!\left(\frac12,n\right).
\]
Then the same-color huge projection sizes satisfy
\[
\sum_{x\in H_I\cap V_R}|\pi_R(X_x(I))|
+
\sum_{x\in H_I\cap V_B}|\pi_B(X_x(I))|
\le \frac{\varepsilon_1}{10}k.
\]
\end{lemma}
\begin{proof}
Let
\[
H_I=\{x:\noderef{TwoBiteHugeClass}(\omega,I,x)\}.
\]
This is the full huge filter used by \(\noderef{HugeFamilySizeBound}\).  The
two same-color filters in the Lean statement are the finite subfilters
\[
H_R=\{x:\noderef{TwoBiteHugeClass}(\omega,I,x)\ \text{and}\ 
      \exists r,\ x=\operatorname{inl}(r)\},
\]
and
\[
H_B=\{x:\noderef{TwoBiteHugeClass}(\omega,I,x)\ \text{and}\ 
      \exists b,\ x=\operatorname{inr}(b)\}.
\]
Equivalently, \(H_R=H_I\cap V_R\) and \(H_B=H_I\cap V_B\), but the displayed
form is the exact filtered-universe form appearing in Lean.  Thus \(x\in H_R\)
means both \(\noderef{TwoBiteHugeClass}(\omega,I,x)\) and
\(x=\operatorname{inl}(r)\) for some unique red base vertex \(r\); similarly
\(x\in H_B\) means \(x=\operatorname{inr}(b)\) for a unique blue base vertex
\(b\).  All cardinalities below are natural finite-set cardinalities before
being cast to \(\mathbb R\), exactly as in the Lean statement.
The displayed equalities \(m=\noderef{TwoBiteNaturalReducedVertexCount}(n)\)
and \(p=\noderef{TwoBiteEdgeProbability}(\frac12,n)\) identify the base size
and edge probability of the sample with the rounded \(m\) and \(p\) used in
\(\noderef{ParameterHierarchyEventualComparisons}\).  Every later citation of
a hierarchy comparison involving \(pm\) uses this identification.

First fix \(x\in H_R\).  Then \(x=\operatorname{inl}(r)\) for a red base
vertex \(r\).  If \(z\in X_x(I)\), then unfolding \(\noderef{TwoBiteX}\) gives
\[
z\in I
\quad\text{and}\quad
z\text{ lies in the lifted neighborhood of }\operatorname{inl}(r).
\]
Unfolding the red branch of the lifted-neighborhood definition recorded in
\(\noderef{FiberAndDegreeEvent}\), the second condition is exactly the red
base-edge condition
\[
G_R.Adj(r,\noderef{TwoBiteRedProjection}(\omega,z)).
\]
In Lean this lifted-neighborhood branch is the filtered-universe condition
\[
(G_R).Adj\ r\ (\noderef{TwoBiteRedProjection}(\omega,z)).
\]
Let
\[
N_R(r)=\{u:\ G_R.Adj(r,u)\}
\]
be the finite red base-neighbor set.  Its cardinality is exactly the finite
degree count used in the red base-degree clause of
\(\noderef{FiberAndDegreeEvent}\): in the Lean definitions this neighbor set
is
\[
\mathrm{univ}.filter\bigl(u\mapsto G_R.Adj(r,u)\bigr),
\]
so its cardinality is
\[
\mathrm{GraphDegreeCount}(G_R,r).
\]
Thus the red base-degree conjunct of \(\noderef{FiberAndDegreeEvent}\) applies
to precisely the real cardinality \((|N_R(r)|:\mathbb R)\), with target
\(pm\).  Unfolding the relative-error predicate in that conjunct gives
\[
\bigl||N_R(r)|-pm\bigr|\le\varepsilon_2pm.
\]
Now take an arbitrary element
\[
u\in(\noderef{TwoBiteX}(\omega,I,x)).\mathrm{image}
      (\noderef{TwoBiteRedProjection}(\omega)).
\]
By the definition of finite-set image, there is some \(z\in X_x(I)\) with
\(u=\noderef{TwoBiteRedProjection}(\omega,z)\).  The preceding paragraph gives
\(G_R.Adj(r,\noderef{TwoBiteRedProjection}(\omega,z))\), and substituting
\(u\) gives \(u\in N_R(r)\).  Thus the exact image appearing in the Lean
statement is contained in that neighbor set:
\[
(\noderef{TwoBiteX}(\omega,I,x)).\mathrm{image}
      (\noderef{TwoBiteRedProjection}(\omega))
  \subseteq N_R(r). \tag{1}
\]
Both sets are finite subsets of the red base vertex set, so monotonicity of
finite cardinality applied to (1), followed by the real cast used in the Lean
statement, gives
\[
|(\noderef{TwoBiteX}(\omega,I,x)).\mathrm{image}
      (\noderef{TwoBiteRedProjection}(\omega))|
\le |N_R(r)|. \tag{2}
\]
The upper side of this absolute-value inequality is
\[
|N_R(r)|-pm\le\varepsilon_2pm,
\]
and adding \(pm\) to both sides gives
\[
|N_R(r)|\le pm+\varepsilon_2pm=(1+\varepsilon_2)pm.
\]
The conjunct
\[
0\le\varepsilon_2\le1
\]
of \(\noderef{ParameterHierarchyEventualComparisons}\), together with the
same package's comparison \(1\le2pm\) for the rounded parameters just
identified, gives \(0\le pm\).  Since \(0\le pm\) and \(\varepsilon_2\le1\),
multiplication by \(pm\) gives \(\varepsilon_2pm\le pm\), and hence
\[
(1+\varepsilon_2)pm\le 2pm.
\]
Combining this with (2), every \(x\in H_R\) satisfies
\[
|(\noderef{TwoBiteX}(\omega,I,x)).\mathrm{image}
      (\noderef{TwoBiteRedProjection}(\omega))|
\le 2pm. \tag{3}
\]

The blue case has the same steps with colors exchanged.  Fix
\(x\in H_B\), so \(x=\operatorname{inr}(b)\) for some blue base vertex \(b\).
If \(z\in X_x(I)\), then unfolding \(\noderef{TwoBiteX}\) and the blue branch
of the lifted-neighborhood definition recorded in
\(\noderef{FiberAndDegreeEvent}\) gives
\[
G_B.Adj(b,\noderef{TwoBiteBlueProjection}(\omega,z)).
\]
Let
\[
N_B(b)=\{u:\ G_B.Adj(b,u)\}.
\]
As in the red case, this is the filtered-universe neighbor set whose
cardinality is the exact finite degree count
\[
\mathrm{GraphDegreeCount}(G_B,b)
\]
appearing in the blue base-degree conjunct of
\(\noderef{FiberAndDegreeEvent}\).  Hence, if
\[
u\in(\noderef{TwoBiteX}(\omega,I,x)).\mathrm{image}
      (\noderef{TwoBiteBlueProjection}(\omega)),
\]
then \(u=\noderef{TwoBiteBlueProjection}(\omega,z)\) for some \(z\in X_x(I)\),
and applying the displayed adjacency to this \(z\) gives \(u\in N_B(b)\).
Therefore the exact image
\((\noderef{TwoBiteX}(\omega,I,x)).\mathrm{image}
      (\noderef{TwoBiteBlueProjection}(\omega))\) is contained in \(N_B(b)\).
Taking finite cardinalities and casting to \(\mathbb R\) gives an upper bound
by \(|N_B(b)|\).  Unfolding the relative-error predicate in the blue
base-degree conjunct of \(\noderef{FiberAndDegreeEvent}\), again with target
\(pm\), gives
\[
\bigl||N_B(b)|-pm\bigr|\le\varepsilon_2pm,
\]
and the same upper-side absolute-value calculation, \(0\le\varepsilon_2\le1\),
and \(1\le2pm\) then give
\[
|(\noderef{TwoBiteX}(\omega,I,x)).\mathrm{image}
      (\noderef{TwoBiteBlueProjection}(\omega))|
\le2pm \qquad (x\in H_B). \tag{4}
\]

In the remaining displayed sums, \(|\pi_R(X_x(I))|\) and
\(|\pi_B(X_x(I))|\) are shorthand for the real cardinalities of the exact
finite images appearing in (3) and (4), respectively.
Summing (3) and (4) over the two finite filters yields
\[
\sum_{x\in H_R}|\pi_R(X_x(I))|
+
\sum_{x\in H_B}|\pi_B(X_x(I))|
\le (|H_R|+|H_B|)\,2pm. \tag{5}
\]
This is a real inequality.  For the red filter, applying finite-sum
monotonicity to the pointwise bound (3) gives
\[
\sum_{x\in H_R}|\pi_R(X_x(I))|
\le (|H_R|:\mathbb R)\,2pm,
\]
because the sum of the constant function \(2pm\) over \(H_R\) is the real cast
of \(|H_R|\) multiplied by \(2pm\).  The same argument with (4) gives the blue
bound, and adding the two inequalities gives (5).

The filters \(H_R\) and \(H_B\) are disjoint.  Indeed, if
\(x\in H_R\cap H_B\), then for some \(r,b\) one would have simultaneously
\(x=\operatorname{inl}(r)\) and \(x=\operatorname{inr}(b)\), impossible for the
sum type of red and blue base vertices.  Moreover both filters are subfamilies
of the full huge filter
\[
H_I=\{x:\noderef{TwoBiteHugeClass}(\omega,I,x)\}.
\]
Consequently their union is a subfamily of \(H_I\).  Since the filters are
finite and disjoint, the natural-cardinality identity for a disjoint union
gives
\[
|H_R\cup H_B|=|H_R|+|H_B|.
\]
Monotonicity of cardinality for the inclusion \(H_R\cup H_B\subseteq H_I\)
then gives
\[
|H_R|+|H_B|=|H_R\cup H_B|\le |H_I|. \tag{6}
\]
After casting these natural cardinalities to \(\mathbb R\), (6) is the real
inequality
\[
(|H_R|:\mathbb R)+(|H_B|:\mathbb R)\le (|H_I|:\mathbb R).
\]
Using \(2pm\ge0\), which follows from \(1\le2pm\), multiplication preserves
this inequality; together with (5) this gives
\[
\sum_{x\in H_R}|\pi_R(X_x(I))|
+
\sum_{x\in H_B}|\pi_B(X_x(I))|
\le |H_I|\,2pm. \tag{7}
\]

Apply \(\noderef{HugeFamilySizeBound}\) with the same
\(\omega,I,k,n,\varepsilon,\varepsilon_1,\varepsilon_2,n_0\), the same
\(\noderef{ParameterHierarchyEventualComparisons}\) hypothesis, the same
strict inequality \(n_0<n\), the same equality
\[
k=\noderef{TwoBiteNaturalIndependenceScale}(1+\varepsilon,n),
\]
the same cardinality hypothesis \(|I|=k\), and the same
\(\noderef{FiberAndDegreeEvent}(\omega,\varepsilon_2)\).  These are exactly
the hypotheses required by \(\noderef{HugeFamilySizeBound}\); the additional
rounded-parameter hypotheses on \(m\) and \(p\) are used only for the
hierarchy comparisons above and below.  Its full huge filter is precisely
\[
H_I=\{x:\noderef{TwoBiteHugeClass}(\omega,I,x)\},
\]
and its second conclusion gives the real-cardinality estimate
\[
|H_I|\le\frac{2k}{t_1}. \tag{8}
\]
Multiplying (8) by the nonnegative real \(2pm\) gives
\[
|H_I|\,2pm
\le
\left(\frac{2k}{t_1}\right)(2pm). \tag{9}
\]
The same-color projection-size comparison in
\(\noderef{ParameterHierarchyEventualComparisons}\), at the present
\(n>n_0\) and for the identified rounded \(m\) and \(p\), is exactly
\[
\left(\frac{2k}{t_1}\right)(2pm)
\le
\frac{\varepsilon_1}{10}k. \tag{10}
\]
Combining (7), (9), and (10) proves the displayed same-color projection-size
bound.
\end{proof}
